% =====================================================================
%  Laboratories for Universal Negatives --- draft 1
%  L. G. Meredith, F1R3FLY.io
% =====================================================================
\documentclass[11pt]{article}

\usepackage[utf8]{inputenc}
\usepackage{amsmath,amssymb,amsthm}
\usepackage[dvipsnames]{xcolor}
\usepackage[margin=1.1in]{geometry}
\usepackage{mathrsfs}
\usepackage{enumitem}
\setlist[itemize]{label=$\bullet$}
\usepackage{booktabs}
\usepackage[expansion=false]{microtype}
\usepackage[hidelinks]{hyperref}

\theoremstyle{definition}
\newtheorem{definition}{Definition}[section]
\newtheorem{construction}[definition]{Construction}
\newtheorem{notation}[definition]{Notation}
\newtheorem{hypothesis}[definition]{Hypothesis}
\theoremstyle{plain}
\newtheorem{proposition}[definition]{Proposition}
\newtheorem{lemma}[definition]{Lemma}
\newtheorem{theorem}[definition]{Theorem}
\newtheorem{corollary}[definition]{Corollary}
\newtheorem{conjecture}[definition]{Conjecture}
\theoremstyle{remark}
\newtheorem{remark}[definition]{Remark}
\newtheorem{example}[definition]{Example}

% --- rho notation, as in "Namespaces Without a Name Server" ---
\newcommand{\rhoc}{\mathrm{rho}}
\newcommand{\quo}[1]{@#1}
\newcommand{\drp}[1]{{*}#1}
\newcommand{\dr}[1]{\overline{#1}}
\newcommand{\inp}[3]{\mathsf{for}(\,#1 \leftarrow #2\,)#3}
\newcommand{\out}[2]{#1!(#2)}
\newcommand{\enc}[1]{[\![#1]\!]}
\newcommand{\nequiv}{\equiv_{\mathcal{N}}}
\newcommand{\nil}{\mathbf{0}}
\newcommand{\Ctx}{\mathcal{C}}
\newcommand{\Obs}{\mathcal{D}}
\newcommand{\bisim}{\approx}
\newcommand{\sbisim}{\sim}
\newcommand{\treq}{\equiv_{\mathrm{tr}}}
\newcommand{\fn}{\mathrm{fn}}
\newcommand{\wild}{\mathord{\underline{\hspace{0.5em}}}}
\DeclareRobustCommand{\S}{\ifmmode\mathsection\else$\mathsection$\fi}
\newcommand{\hole}{[\,\cdot\,]}
\newcommand{\ok}{\mathsf{ok}}
\newcommand{\go}{\mathsf{go}}
\newcommand{\lk}{\mathsf{l}}
\newcommand{\Lab}[1]{\xrightarrow{\,#1\,}}
\newcommand{\WLab}[1]{\overset{#1}{\Longrightarrow}}
\newcommand{\act}{\mathbin{\star}}
\newcommand{\cmp}{\mathbin{\circ}}
\newcommand{\Dup}{\mathsf{D}}
\newcommand{\ABS}{\textup{\textsc{Absence}}}
\newcommand{\TREQ}{\textup{\textsc{TraceEq}}}
\newcommand{\Pclass}{\mathbf{P}}
\newcommand{\NPclass}{\mathbf{NP}}
\newcommand{\coNP}{\mathbf{coNP}}
\newcommand{\cost}{\mathrm{cost}}
\newcommand{\Reg}{\mathcal{R}}

\title{\textbf{Laboratories for Universal Negatives}\\[4pt]
\large $\Pclass$ versus $\NPclass$ as the existence of an experimental regime,\\
in the rho calculus with context-labelled bisimilarity}
\author{L.\,G.\ Meredith\\[2pt] \normalsize F1R3FLY.io\quad \texttt{lgreg.meredith@f1r3fly.io}}
\date{October 2026 --- draft 1}

\begin{document}
\maketitle

\begin{abstract}
\noindent
We recast the separation of $\Pclass$ from $\NPclass$, if there is one,
as a statement about which experiments a process can perform on another
process. The setting is the rho calculus with context-labelled
transitions: a label is a context, and an observer class is a set of
contexts. Three steps. First, on bounded processes bisimilarity is
polynomial and trace equivalence is $\coNP$-complete, and $\Pclass =
\NPclass$ holds exactly when a polynomial-time normaliser turns trace
equivalence into bisimilarity; a faithful normaliser whose output is
determinate provably does not exist. Second, internalising the normaliser as a rho
process $\enc{\Phi}(x)$ that receives its input as code on $x$ shows that
everything it can do with that code is interaction: the context
$\enc{\Phi}(x) \mid \out{x}{\hole}$ places its hole under a lift, where
communication depends on name equality, and the only uses of a received
name are as a channel and under a drop. What distinguishes code from a
live process is that code supports repeated experiments. Third, we
sharpen the graded interaction of fuzzy rho into an apparatus for
\emph{controlled} experiments --- stamping for isolation and attribution,
positional interventions, a unit-grade control, pooled readout --- and
define an experimental regime as an access discipline, an intervention
language, a readout and a cost model. The result is a classification.
Black-box regimes are provably unable to certify that a bounded process
cannot succeed: an adversary argument costs them $2^d$ intervened copies
on inputs of size $O(d \log d)$. Under two stated cost hypotheses, $\Pclass = \NPclass$ if
and only if some admissible regime certifies such universal negatives at
polynomial cost. The black-box separation is an in-calculus analog of
an oracle separation, so what remains of $\Pclass \neq \NPclass$ in this
language is entirely a statement about white-box access.
\end{abstract}

% =====================================================================
\section{Introduction}\label{sec:intro}
% =====================================================================

This note began as a discussion of whether $\Pclass \neq \NPclass$ could
be independent of ZFC, and of whether forcing could be adapted to show it.
The obstruction to the second question is well understood. The statement
$\Pclass \neq \NPclass$ is arithmetical ($\Pi^0_2$). Forcing adds no
new natural numbers, so arithmetical statements are absolute between a
model and its forcing extensions, and no set-theoretic forcing changes
the truth value of $\Pclass \neq \NPclass$. (Shoenfield's theorem
\cite{Shoenfield61} carries absoluteness much further, to $\Sigma^1_2$.)
Forcing survives only in its model-theoretic form, over nonstandard
models of weak arithmetic \cite{Krajicek11}. Related unprovability
results for circuit lower bounds in bounded arithmetic are obtained by
natural-proofs arguments rather than by forcing \cite{Razborov95}.
Independence from strong theories would also be close to a positive
answer. Ben-David and Halevi \cite{BenDavidHalevi92} show that if
$\Pclass \neq \NPclass$ is independent of PA augmented with all true
$\Pi_1$ sentences --- a strong, non-recursively-enumerable theory ---
then $\NPclass$ has almost-polynomial algorithms. See \cite{Aaronson03}
for a survey.

Rather than pursue independence directly, we ask what shape the
obstruction takes when the question is posed in a calculus whose native
notion of equivalence is behavioural. The answer turns out to be
operational: it is a question about laboratories.

\subsection*{The idea, in four lines}

\begin{enumerate}[label=(\arabic*),leftmargin=2.2em]
\item Branching time is cheap and linear time is not: on bounded
processes, bisimilarity is decidable in polynomial time and trace
equivalence is $\coNP$-complete (\S\ref{sec:bounded}).

\item $\Pclass = \NPclass$ if and only if a polynomial-time normaliser
sends trace equivalence to bisimilarity (\S\ref{sec:normaliser}).

\item Internalise the normaliser as a rho process $\enc{\Phi}(x)$
receiving its input as code on $x$. Everything it does with that code is
an experiment, and code differs from a live process in admitting repeated
experiments (\S\ref{sec:internal}--\S\ref{sec:tiers}).

\item $\Pclass \neq \NPclass$ if and only if no admissible experimental
regime lets $\enc{\Phi}(x)$ certify a universal negative at polynomial
cost. Black-box regimes provably cannot; white-box regimes are exactly as
hard as $\Pclass$ versus $\NPclass$ (\S\ref{sec:experiments}--\S\ref{sec:regimes}).
\end{enumerate}

\subsection*{Contributions}

\begin{enumerate}[label=(\roman*),leftmargin=2.2em]
\item A bisimulation-native statement of $\Pclass$ versus $\NPclass$
(Theorem~\ref{thm:normaliser}), with an unconditional lower bound for
faithful determinate normalisers (Proposition~\ref{prop:determinate}).

\item The observation that a reflective normaliser is a context whose hole
lies under a lift, and that in such a context every operation on the
input is an interaction (Lemma~\ref{lem:uses}); the access tiers that
follow are distinguished by the structural rules they permit
(\S\ref{sec:tiers}).

\item An apparatus for controlled experiments built from the graded
interaction of fuzzy rho \cite{FuzzyRho}, with the two conditions ---
inert signatures and equivariance --- under which a stamped copy is a
faithful copy, and a counterexample showing that flat stamping is not
equivariant (Example~\ref{ex:nonequivariant}).

\item A definition of experimental regime, an exponential lower bound for
black-box regimes that is unconditional in the number of experiments
(Theorem~\ref{thm:blackbox}), and a
characterisation of $\Pclass = \NPclass$ by the existence of an
admissible regime (Theorem~\ref{thm:whitebox}), relative to two stated
cost hypotheses.

\item A heuristic map from the known barriers to features of the
apparatus (\S\ref{sec:barriers}).
\end{enumerate}

\noindent Appendix~\ref{app:status} records the status of each result.
Several are proved; some are proved relative to stated hypotheses; some
are sketched; some are conjectures, and are labelled as such. Nothing
here claims progress on $\Pclass$ versus $\NPclass$ itself. The
contribution is a reformulation whose parts are the parts of a
laboratory, together with the place in that laboratory where the known
barriers sit.

% =====================================================================
\section{The rho calculus and context-labelled bisimilarity}\label{sec:rho}
% =====================================================================

\begin{definition}[Syntax]
Processes $P$ and names $x$ are given by mutual recursion:
\[
P \ ::=\ \nil \ \mid\ P_1 \mid P_2 \ \mid\ \out{x}{P}
   \ \mid\ \mathsf{for}(y_1 \leftarrow x_1 \,\&\, \cdots \,\&\, y_n \leftarrow x_n)P
   \ \mid\ \drp{x}
\qquad\qquad
x \ ::=\ \quo{P}
\]
$\out{x}{P}$ is the \emph{lift}: it quotes $P$ and outputs the name
$\quo{P}$ on $x$. The input binds $y_1,\dots,y_n$ in $P$; for $n = 1$ we
write $\inp{y}{x}{P}$, and for $n > 1$ it is a \emph{join}. $\drp{x}$ is
the \emph{drop}. We write $\wild$ for a bound name that does not occur.
A process is \emph{closed} if no variable bound by an input occurs free in
it; every quoted name in a closed process is then a constant.
\end{definition}

This is the calculus of \cite{MeredithRadestock05} with joins, in the
notation of \cite{NWNS}. A name is in \emph{normal form} if no quote in it
immediately encloses a drop. Rewriting $\quo{(\drp{x})}$ to $x$
throughout gives a normal form $\nequiv$-equal to any given name. Name
equivalence is decidable in polynomial time, being equality of unordered
trees after normalisation.

\begin{definition}[Name equivalence, structural congruence, reduction]
As in \cite{NWNS}: $\nequiv$ is the least equivalence on names with
$\quo{P_1} \nequiv \quo{P_2}$ whenever $P_1 \equiv P_2$, and
$\quo{(\drp{x_1})} \nequiv x_2$ whenever $x_1 \nequiv x_2$; $\equiv$ is
the least congruence containing $\alpha$-equivalence and the abelian
monoid laws for $\mid$ with unit $\nil$, names being compared up to
$\nequiv$. Reduction is
\[
\frac{x_i \nequiv x_i' \quad (1 \le i \le n)}
     {\mathsf{for}(y_1 \leftarrow x_1 \,\&\, \cdots \,\&\, y_n \leftarrow x_n)P
      \ \mid\ \out{x_1'}{Q_1} \mid \cdots \mid \out{x_n'}{Q_n}
      \ \to\ P\{\quo{Q_1}/y_1, \dots, \quo{Q_n}/y_n\}}
\]
closed under $\mid$ and $\equiv$. Substitution replaces all names
$\nequiv$-equal to the variable, reaches the object of a lift, does not
reach beneath a quote, and performs the drop:
$(\drp{x})\{\quo{P}/y\} = P$ when $x \nequiv y$.
\end{definition}

The duplicator $\Dup(x) \triangleq \inp{y}{x}{(\drp{y} \mid
\out{x}{\drp{y}})}$ satisfies $\out{x}{Q} \mid \Dup(x) \to Q \mid
\out{x}{Q}$; it is the source of contraction on names.

\subsection{Context-labelled transitions}

We do not use barbs. Following \cite{NWNS}, observation is by contexts.

\begin{definition}[Context-labelled transitions and bisimilarity]
\label{def:clbisim}
Write $P \Lab{C} P'$ when $C[P] \to P'$, and $P \WLab{C} P'$ when $C[P]
\to^* P'$. For an \emph{observer class} $\Obs \subseteq \Ctx$, a symmetric
relation $\mathcal{R}$ is a $\Obs$-bisimulation if $P \mathrel{\mathcal{R}}
Q$ and $P \Lab{C} P'$ with $C \in \Obs$ imply $Q \WLab{C} Q'$ with $P'
\mathrel{\mathcal{R}} Q'$; $\bisim_\Obs$ is the largest such.
\end{definition}

\begin{definition}[Traces]
An $\Obs$-trace of $P$ is a sequence $C_1 \cdots C_k$ of contexts in
$\Obs$ with $P \WLab{C_1} P_1 \WLab{C_2} \cdots \WLab{C_k} P_k$. $P
\treq^\Obs Q$ when $P$ and $Q$ have the same $\Obs$-traces. Plainly
$\bisim_\Obs \;\subseteq\; \treq^\Obs$.
\end{definition}

\begin{definition}[Success]\label{def:success}
Fix a name $\ok$ and let $C_\ok = \hole \mid \inp{\wild}{\ok}{\nil}$.
$P$ \emph{can succeed} if $P \WLab{\hole} P' \Lab{C_\ok} P''$ for some
$P', P''$ in which the displayed step consumes an output on $\ok$.
\end{definition}

Success is thus a context-labelled transition, not a barb. The whole of
this note is about the choice of $\Obs$: an experimental regime will be,
among other things, an observer class.

\begin{remark}[Holes under a lift]\label{rem:bihol}
The companion paper \cite{BIHOL} shows that in rho, context-labelled
bisimilarity is a congruence for parallel composition, lift, input and
drop-of-a-quote when observations are suitably restricted, and fails to
be one for contexts with the hole under a quote, because communication
depends on name equality rather than on bisimilarity. The contexts that
matter in \S\ref{sec:internal} are exactly of that kind.
\end{remark}

\subsection{Cost}

Cost-accounted rho charges phlogiston for communication. We need only two
properties of the charge, stated as hypotheses because they are
properties of a cost model, not of the calculus.

\begin{hypothesis}[Simulation]\label{hyp:sim}
There is a polynomial $q$ such that, for every deterministic Turing
machine $M$, there is a rho process $R_M(x)$ which, given an input word
$w$ as data on $x$, runs determinately to a representation of $M(w)$ at
phlogiston at most $q(t_M(w) + |w|)$, where $t_M(w)$ is the running time.
\end{hypothesis}

\begin{hypothesis}[Accounting]\label{hyp:acct}
Each communication --- including the graded communications of
\S\ref{sec:experiments} and the probe interactions of
\S\ref{sec:tiers} --- is charged at least $1$ plus the total size of the
transmitted terms. Consequently a run of phlogiston $c$ from a term of
size $n$ can be simulated by a deterministic Turing machine in time
polynomial in $c + n$.
\end{hypothesis}

Hypothesis~\ref{hyp:acct} excludes the size explosion that free
duplication would otherwise allow: a process that repeatedly receives $Q$
and sends $Q \mid Q$ pays for every copy it makes.

% =====================================================================
\section{Bounded processes}\label{sec:bounded}
% =====================================================================

Two presentations of bounded behaviour are needed: an explicit one, on
which the bisimulation analog is stated, and a succinct one, on which
experiments are performed. They are polynomially inter-reducible through
$\coNP$-completeness.

\subsection{Explicit presentation}

\begin{definition}
An \emph{explicit bounded process} is a finite acyclic labelled
transition system $G = (S, \Sigma, \to, s_0)$, given as a list of
transitions; $|G|$ is the length of that list. $G$ is \emph{determinate}
if every state has at most one $a$-successor for each $a \in \Sigma$.
$\sbisim$ is strong bisimilarity and $\treq$ trace equivalence.
\end{definition}

\begin{construction}[Encoding]\label{con:encoding}
Fix pairwise distinct closed names $\go_s, \lk_s$ for $s \in S$ and $a$
for $a \in \Sigma$, chosen injectively (for instance by quoting numerals).
Put
\[
\begin{aligned}
\mathcal{E}(s) &\;=\; \out{\lk_s}{\nil}
  \ \mid\ \prod_{(s,a,t) \in \to}
     \mathsf{for}(\wild \leftarrow a \,\&\, \wild \leftarrow \lk_s)\,
     \out{\go_t}{\nil},\\
\mathcal{E}(G) &\;=\; \out{\go_{s_0}}{\nil}
  \ \mid\ \prod_{s \in S} \inp{\wild}{\go_s}{\mathcal{E}(s)}.
\end{aligned}
\]
The lock $\lk_s$ lets exactly one outgoing transition of $s$ fire; a
state shared by several paths is encoded once and triggered at most once
per run, because $G$ is acyclic. $|\mathcal{E}(G)| = O(|G|)$. The
observer class is $\Obs_\Sigma = \{\, \hole \mid \out{a}{\nil} : a \in
\Sigma \,\}$.
\end{construction}

Communications on $\go$-names are \emph{administrative}; a join
consuming an offered $\out{a}{\nil}$ is \emph{principal}, in the sense of
the phase separation of \cite{NWNS}. Administrative reductions are
deterministic (each $\go_t$ has one listener) and terminate. A losing
branch of $s$ is left as a join on $\lk_s$, which is never sent again
and which no context in $\Obs_\Sigma$ can name, so it is inert.

\begin{proposition}[Adequacy of the encoding]\label{prop:adequacy}
Taking transitions to be principal steps followed by administrative
normalisation, $s \sbisim t$ iff $\mathcal{E}(s) \bisim_{\Obs_\Sigma}
\mathcal{E}(t)$, and $s \treq t$ iff $\mathcal{E}(s)
\treq^{\Obs_\Sigma} \mathcal{E}(t)$.
\end{proposition}

\begin{proof}[Proof sketch]
After administrative normalisation, a term reachable from
$\mathcal{E}(G)$ is $\mathcal{E}(s) \mid R$ with $R$ a parallel
composition of inert joins and spent triggers. A principal step on
$\out{a}{\nil}$ is available from it exactly when $s \Lab{a} t$ for some
$t$, and leads, after normalisation, to $\mathcal{E}(t) \mid R'$. The
residue $R$ contributes no $\Obs_\Sigma$-transitions. Both equivalences
then transfer along the correspondence.
\end{proof}

\begin{remark}
The restriction to $\Obs_\Sigma$ matters. An arbitrary rho context can
name $\lk_s$, since it is computed by the encoding, and so revive a
losing branch, much as an arbitrary context separates encodings in the
necessity result of \cite{NWNS}. Reviving a branch uses knowledge of the
term, which is a white-box operation in the sense of \S\ref{sec:tiers}.
\end{remark}

\begin{theorem}[Known complexity]\label{thm:known}
On explicit bounded processes:
\begin{enumerate}[label=(\alph*),leftmargin=2.2em,nosep]
\item $\sbisim$ is decidable in polynomial time \cite{PaigeTarjan87};
on finite transition systems in general it is $\Pclass$-complete
\cite{BGS92};
\item $\treq$ is $\coNP$-complete.
\end{enumerate}
\end{theorem}

\begin{proof}[Proof of (b), sketch]
Inequivalence is in $\NPclass$: guess a trace, of length at most the
depth, present in one system and absent from the other; membership of a
given trace is decided by tracking the set of reachable states.
Hardness: equivalence of regular expressions without star is
$\coNP$-hard \cite{StockmeyerMeyer73}; such an expression converts in
polynomial time to an acyclic automaton, and marking acceptance by a
final $\ok$-transition and trimming states that cannot reach acceptance
turns language equivalence into trace equivalence.
\end{proof}

\subsection{Succinct presentation}

\begin{definition}\label{def:flat}
A process is \emph{flat} if it is closed, contains no drop (not even beneath a quote), and no bound
name occurs in subject position. In a flat process each communication consumes an input
prefix not under a lift, and none is created, since no transmitted term is
ever run. So every run of a flat $P$ has length at most $|P|$. Let
$\ABS$ be the problem: given a flat $P$, decide that $P$ \emph{cannot}
succeed.
\end{definition}

\begin{proposition}\label{prop:abscomplete}
$\ABS$ is $\coNP$-complete.
\end{proposition}

\begin{proof}
It is in $\coNP$, since a successful run has length at most $|P|$ and is a
certificate of success, checkable in polynomial time because name
equivalence is. For hardness, reduce from CNF unsatisfiability.
Given variables $v_1, \dots, v_d$ and clauses $c_1, \dots, c_m$, fix
pairwise distinct closed names $b_i, t_{ij}, f_{ij}, e_j, \ok$, and let
\[
\mathit{Race}_i \;=\; \out{b_i}{\nil}
  \ \mid\ \inp{\wild}{b_i}{\textstyle\prod_{j : v_i \in c_j} \out{t_{ij}}{\nil}}
  \ \mid\ \inp{\wild}{b_i}{\textstyle\prod_{j : \neg v_i \in c_j} \out{f_{ij}}{\nil}},
\]
let $\mathit{Clause}_j$ be the parallel composition of
$\inp{\wild}{t_{ij}}{\out{e_j}{\nil}}$ for each positive occurrence of
$v_i$ in $c_j$ and $\inp{\wild}{f_{ij}}{\out{e_j}{\nil}}$ for each
negative one, and put
\[
P_\varphi \;=\; \prod_i \mathit{Race}_i \ \mid\ \prod_j \mathit{Clause}_j
  \ \mid\ \mathsf{for}(\wild \leftarrow e_1 \,\&\, \cdots \,\&\,
  \wild \leftarrow e_m)\, \out{\ok}{\nil}.
\]
Each race is won by exactly one listener, which fixes a truth value for
$v_i$; a clause signals $e_j$ only if one of its literals was made true;
the final join fires only if every clause signalled. A satisfying
assignment yields a successful schedule and conversely, so $P_\varphi$
can succeed iff $\varphi$ is satisfiable. $P_\varphi$ is flat and of size
polynomial in $\varphi$.
\end{proof}

The races are where scheduling nondeterminism lives. A \emph{resolution}
of a flat process assigns a winner to each race. The point of
\S\ref{sec:experiments} is to make resolutions the experimenter's choice
rather than the scheduler's.

% =====================================================================
\section{Normalisers}\label{sec:normaliser}
% =====================================================================

\begin{definition}\label{def:normaliser}
A \emph{normaliser} is a polynomial-time map $\Phi$ on explicit bounded
processes such that for all $G, H$
\[
G \treq H \quad\Longleftrightarrow\quad \Phi(G) \sbisim \Phi(H).
\]
It is \emph{faithful} if moreover $\Phi(G) \treq G$, and \emph{determinate}
if every $\Phi(G)$ is determinate.
\end{definition}

\begin{theorem}\label{thm:normaliser}
$\Pclass = \NPclass$ if and only if a normaliser exists, and then a
faithful one exists.
\end{theorem}

\begin{proof}
($\Leftarrow$) Given $\Phi$, decide $G \treq H$ by computing $\Phi(G)$ and
$\Phi(H)$ and testing $\sbisim$ (Theorem~\ref{thm:known}(a)). This puts a
$\coNP$-complete problem in $\Pclass$, so $\coNP = \Pclass$ and hence
$\NPclass = \Pclass$.

($\Rightarrow$) If $\Pclass = \NPclass$ the polynomial hierarchy collapses
to $\Pclass$. Let $m(G)$ be the least size of an explicit process
trace-equivalent to $G$, and let $\Phi(G)$ be the lexicographically least
explicit process $G'$, encoded as a list of transitions over states
numbered $0, \dots, m$ and over the labels occurring in traces of $G$, with $|G'| = m(G)$ and
$G' \treq G$. Both depend only on the trace class of $G$. Since $m(G) \le
|G|$, $m(G)$ is found by asking, for each $m \le |G|$, whether some $G'$
of size $m$ is trace-equivalent to $G$. The bits of $\Phi(G)$ are then
found one at a time by asking whether some such $G'$ extends a given
prefix. Each question has the form: there is a $G'$ such that for all
traces up to the depth bound, membership agrees. That is
$\Sigma^p_2$, hence polynomial-time. (This is the standard fact that
$\Pclass = \NPclass$ yields polynomial-time canonical forms for
$\coNP$ equivalence relations.) Then $\Phi(G) \treq G$, so
$\Phi$ is faithful. If $G \treq H$ then $\Phi(G) = \Phi(H)$, hence
$\Phi(G) \sbisim \Phi(H)$. Conversely $\Phi(G) \sbisim \Phi(H)$ implies
$\Phi(G) \treq \Phi(H)$, hence $G \treq H$.
\end{proof}

\begin{proposition}[Faithful determinate normalisers do not exist]\label{prop:determinate}
There are explicit bounded processes $G_n$ with $|G_n| = O(n^2)$ such
that every determinate $H$ with $H \treq G_n$ has at least $2^n$ states.
In particular no faithful determinate normaliser runs in polynomial time,
unconditionally.
\end{proposition}

\begin{proof}
Let $L_n = \{\, uv \in \{0,1\}^{2n} : u \neq v,\ |u| = |v| = n \,\}$, and
let $G_n$ have the traces of the words of $L_n$ each followed by $\ok$.
$G_n$ guesses a position $i \le n$, records the bit read there, counts
$n$ further positions, checks that the bit there differs, reads to length
$2n$ and emits $\ok$. Its states are indexed by a position, a guessed
index and a bit, so there are $O(n^2)$ of them, and it is acyclic. Let $H$
be determinate with $H \treq G_n$. Each $u \in \{0,1\}^n$ is a trace of
$H$, since $u u'$ is in $L_n$ for any $u' \ne u$, and so leads to a
unique state $h_u$. For $u \neq u'$ the trace $u' \ok$ is available from
$h_u$ but not from $h_{u'}$, since $u'u' \notin L_n$. So the $h_u$ are
pairwise distinct. This is the classical blow-up of determinisation
\cite{MeyerFischer71,Moore71}, specialised to finite languages.
\end{proof}

Faithfulness matters. Without it a determinate normaliser exists iff
$\Pclass = \NPclass$: if $\Pclass = \NPclass$, map $G$ to the straight-line
system spelling out the bits of the canonical form of
Theorem~\ref{thm:normaliser}. Straight lines are determinate, and two of
them are bisimilar exactly when the canonical forms are equal.

\begin{remark}[The ladder]\label{rem:ladder}
There are three grades of normaliser. A faithful \emph{determinate}
normaliser is excluded unconditionally (Proposition~\ref{prop:determinate}). An
\emph{arbitrary} polynomial-time normaliser exists iff $\Pclass =
\NPclass$ (Theorem~\ref{thm:normaliser}). Between them sit
\emph{compositional} normalisers, which act on the encoded term
$\mathcal{E}(G)$ in a syntax-directed way and commute with its operators,
parallel composition and prefixing. For
these we expect exponential lower bounds of the subset-construction kind
(Conjecture~\ref{conj:compositional}), much as monotone circuits admit
lower bounds that general circuits have so far resisted.
\end{remark}

\begin{conjecture}\label{conj:compositional}
No faithful compositional normaliser runs in polynomial time.
\end{conjecture}

% =====================================================================
\section{Internalising the normaliser}\label{sec:internal}
% =====================================================================

Write the internalised normaliser as $\enc{\Phi}(x)$: a fixed rho process,
independent of its input, with a free name $x$ on which it expects that
input. The configuration is
\[
\enc{\Phi}(x) \ \mid\ \out{x}{P} \;\;\to^*\;\; P'.
\]
In context-labelled terms the experimenter \emph{is} the context
$C_\Phi = \enc{\Phi}(x) \mid \out{x}{\hole}$, and its hole lies under a
lift. By Remark~\ref{rem:bihol} this is the position where
$\bisim$ fails to be a congruence. So the requirement that $P \bisim Q$
imply $C_\Phi[P] \bisim C_\Phi[Q]$, and \emph{a fortiori} the requirement
that $P \treq Q$ imply it, is not free: it is the burden the normaliser
carries. The live alternative, $\enc{\Phi} \mid \hole$, gets invariance
under $\bisim$ for free from congruence and loses faithfulness instead
(Example~\ref{ex:linear}).

\begin{lemma}[What a received name is good for]\label{lem:uses}
In $\inp{y}{x}{R}$, every free occurrence of $y$ in $R$ is either the
subject of an input or output or the argument of a drop. Consequently,
once $\enc{\Phi}(x)$ has received $\quo{P}$, it can
\begin{enumerate}[label=(\roman*),leftmargin=2.2em,nosep]
\item run $P$, by a drop; or
\item use $\quo{P}$ as a channel, which matters only to processes that
also hold a name $\nequiv \quo{P}$, and so amounts to testing $\quo{P}$
for name equality against names it or its environment constructs; or
\item forward it inside a lifted term, which defers (i) or (ii) to the
receiver.
\end{enumerate}
\end{lemma}

\begin{proof}
By the grammar, names occur only as subjects of input and output, as
binders, under drop, and beneath quotes; occurrences beneath a quote are
not free, since substitution does not reach there. Communication on a name fires only when the two
subjects are $\nequiv$-equal.
\end{proof}

Pure rho has no third operation, no deconstructor. Any test on the syntax
of $P$ is an interaction: either running $P$ against a context, or a
communication whose firing condition is a name equality the experimenter
arranged. So the quoted reading reduces, internally, to possibly iterated
interaction with $P$. What separates it from the live reading is
\emph{resource}: interaction with a live $P$ can use $P$ up, while a name
can be dropped as often as the experimenter pays for.

\begin{example}[Why one experiment is not enough]\label{ex:linear}
Let $G = a.(b + c)$ and $H = a.b + a.c$, so $G \treq H$ and $G \not\sbisim
H$. A context interacting with the live $\mathcal{E}(H)$ learns nothing
until it offers $a$, and in $\mathcal{E}(H)$ the very join that consumes
the offer also chooses the branch. Afterwards the continuation can do $b$
or $c$, not both. The losing branch survives only as a dead join on a
lock, which no observer in $\Obs_\Sigma$ can revive. A faithful normal
form must have both traces $ab$ and $ac$, and a single linear experiment
on $\mathcal{E}(H)$ cannot recover the one it did not see. In a calculus
with summation the losing branch is discarded outright.
\end{example}

% =====================================================================
\section{Access tiers and structural rules}\label{sec:tiers}
% =====================================================================

Quotation gives the input the exponential of linear logic: $\quo{P}$
behaves as $!P$, the drop as dereliction, the duplicator as contraction.
Access tiers are distinguished by which of these rules the experimenter
may use on the received name, and by whether it may inspect the term.

\begin{definition}[Tiers]\label{def:tiers}
\begin{enumerate}[label=(\alph*),leftmargin=2.2em]
\item \emph{Linear}: the received name is dropped at most once, as in the
exactly-once input $\mathsf{for}(y \mathrel{\circ\!\!-} x)P$ of the
linear fragment. No contraction.

\item \emph{Replicable black-box}: the received name is used only under
drop, and is forwarded only to receivers that likewise use it only under
drop (\emph{drop-only} usage). The experimenter's contexts mention, of
the copy's names, only those in a declared public interface (here $\ok$
and any declared intervention ports). Contraction is allowed; inspection
and name guessing are not.

\item \emph{White-box}: as (b), and in addition the experimenter may apply
the probes of the observer extension \cite{ObserverExt} to the received
name itself --- to the quoted code, which does not run while it is read.
The probes make each term constructor observable by interaction: an atom $\mathsf{Ask}_C$ meeting $C(t_1,\dots,t_n)$
yields its arguments, and $\mathsf{Get}_{C,i}$ selects the $i$-th. Under
the capability-guarded refinement of that construction, holding $x$ is
what licenses inspection.
\end{enumerate}
\end{definition}

Name guessing --- use (ii) of Lemma~\ref{lem:uses} against names the
experimenter constructs --- is excluded from tier (b) deliberately. It is
a syntactic test, and it belongs with inspection.

\begin{conjecture}[Drop-only contexts are congruent]\label{conj:dropOnly}
If the hole of $C$ lies under a lift and the corresponding name is used
drop-only, then $P \bisim_\Obs Q$ implies $C[P] \bisim_\Obs C[Q]$ for
suitable $\Obs$. That is, drop-only usage makes a quote transparent, and
tier (b) is an admissible observer class in the sense of \cite{BIHOL}.
\end{conjecture}

% =====================================================================
\section{Controlled experiments: sharpening fuzzy rho}\label{sec:experiments}
% =====================================================================

Repetition alone yields only existential evidence. If a copy of $P$
succeeds, $P$ can succeed. If it fails, the scheduler may simply have
resolved a race the other way. A controlled experiment needs isolation,
attribution, fidelity, intervention, a control and a readout. We take
each from the graded interaction of fuzzy rho \cite{FuzzyRho}, in which
the transmitted term is modulated by a grade comparing the two
interaction surfaces:
\[
\inp{y}{x}{R} \ \mid\ \out{u}{Q} \;\to\; R\{\quo{((u \cmp x) \act Q)}/y\},
\qquad
(u \cmp x) \act Q = Q\{\quo{(\dr{u} \mid \dr{x} \mid \dr{z})}/z : z \in \fn(Q)\},
\]
where $\dr{z}$ is the process quoted by $z$. This is the
\emph{accumulator} model; the \emph{attenuator} replaces $\mid$ by the
multiset meet $\wedge$.

\subsection{Isolation and attribution: the stamp}

\begin{definition}[Stamp]
A \emph{signature} is a closed process $S$ having, as a $\mid$-prime
component, an inert input $K_S = \inp{\wild}{k_S}{\nil}$ on an
unforgeable name $k_S$ that occurs nowhere else. The \emph{stamp}
$\theta_S$ sends each name $z$ in normal form to $\quo{(S \mid \dr{z})}$,
and $S \act P$ applies $\theta_S$ to the names of $P$, taken in normal
form.
\end{definition}

\begin{lemma}[Stamps are injective]\label{lem:injective}
$\theta_S(z_1) \nequiv \theta_S(z_2)$ implies $z_1 \nequiv z_2$.
\end{lemma}

\begin{proof}
$S \mid \dr{z}$ has the input $K_S$ as a component, so it is never
$\equiv$ to a lone drop, and only the first clause of $\nequiv$ applies:
$\theta_S(z_1) \nequiv \theta_S(z_2)$ means $S \mid \dr{z_1} \equiv S
\mid \dr{z_2}$. Up to $\equiv$ a process is a finite multiset of
$\mid$-prime components, and the free commutative monoid is cancellative,
so $\dr{z_1} \equiv \dr{z_2}$, and $z_1 \nequiv z_2$ since both are in
normal form.
\end{proof}

\begin{lemma}[Distinct signatures are disjoint]\label{lem:disjoint}
Let $S_1, S_2$ be signatures with $k_{S_1} \neq k_{S_2}$, and $z_1, z_2$
names of an input $P$ that mentions neither $k_{S_i}$. Then
$\theta_{S_1}(z_1) \not\nequiv \theta_{S_2}(z_2)$.
\end{lemma}

\begin{proof}
If $S_1 \mid \dr{z_1} \equiv S_2 \mid \dr{z_2}$, the prime component
$K_{S_1}$ occurs on the right, hence in $S_2$ or in $\dr{z_2}$. Neither
mentions $k_{S_1}$.
\end{proof}

Run $k$ copies with signatures $S_1, \dots, S_k$ whose unforgeable names
$k_{S_i}$ are distinct, as the content-addressed channels of
cost-accounted rho provide. Then by Lemmas~\ref{lem:injective}
and~\ref{lem:disjoint} the copies share no channel, so they cannot
interfere with one another, and every observation carries the signature
of the copy that made it. The experimenter knows each $S_i$, so it can
compute the stamped interface $\theta_{S_i}(\ok)$ on which to listen.

\subsection{Fidelity: inertness and equivariance}

A stamped copy is useful only if it is a faithful copy. Two things can go
wrong.

\paragraph{Signatures run under drop.} If $P$ contains $\drp{z}$ for a
free $z$, the copy contains $\drp{\theta_S(z)}$, which runs $S \mid
\dr{z}$. The signature must therefore be \emph{inert}: built, for
instance, from inputs on unforgeable names that nothing ever sends on.

\paragraph{Names synthesised at runtime escape the stamp.} Flat
substitution stamps the names $P$ has at the moment it is sent, not the
names it computes while it runs.

\begin{example}[Flat stamping is not equivariant]\label{ex:nonequivariant}
Let $P = \out{a}{\nil} \mid \inp{y}{a}{\out{y}{\nil}} \mid
\inp{\wild}{\quo{\nil}}{\out{\ok}{\nil}}$. Then $P$ receives $y :=
\quo{\nil}$, outputs on $\quo{\nil}$, and succeeds. In $S \act P$, for any
signature $S$, the names $a$, $\ok$ and $\quo{\nil}$ are stamped, but
the transmitted name $\quo{\nil}$ is synthesised from the object of the
lift, which has no names, and so is not. The copy outputs on $\quo{\nil}$
while its listener waits on $\theta_S(\quo{\nil}) \nequiv \quo{S}$. The
copy cannot succeed.
\end{example}

\begin{definition}[Experimental grading]\label{def:expgrading}
A grading $S \act (-)$ is \emph{experimental} if
\begin{enumerate}[label=(E\arabic*),leftmargin=2.6em,nosep]
\item its signatures are signatures in the sense above, pairwise
distinct across copies;
\item its signatures are inert;
\item it is \emph{equivariant}: $P \to P'$ implies $S \act P \to S \act
P'$, and every reduction of $S \act P$ is of this form.
\end{enumerate}
\end{definition}

(E3) is the dynamic form of the binder transport law of \cite{FuzzyRho},
$a \act [y]P = [y](a' \act P)$. It says the copy runs \emph{inside} the
stamp, so that every name it ever creates is stamped.

\begin{proposition}\label{prop:fidelity}
If $S \act (-)$ is experimental, the relation $\{(P', S \act P')\}$ is a
bisimulation between the reduction systems of $P$ and of $S \act P$, with
observations on $z$ matched to observations on $\theta_S(z)$. In
particular $S \act P$ can succeed at $\theta_S(\ok)$ iff $P$ can succeed.
\end{proposition}

\begin{proof}
Immediate from (E3), with (E2) ensuring that the dropped signature
contributes no transitions.
\end{proof}

Flat processes (Definition~\ref{def:flat}) synthesise no names that are
ever used as channels and contain no drops, so flat stamping is
experimental on them. Realising (E3) for general processes needs a
namespace-relative evaluation of the copy; this is open
(\S\ref{sec:open}).

\subsection{Intervention, control, dose and readout}

\paragraph{Intervention.} Uniform stamping isolates a copy but does not
steer it. To control a variable we need the analog of Pearl's
$\mathrm{do}(\cdot)$ \cite{Pearl09}: resolve a chosen race to a chosen
branch. Stamping a single \emph{occurrence} does this. In
$\out{z}{A} \mid \inp{y}{z}{B} \mid \inp{y}{z}{C}$, stamping the subject
of the third component with an inert tag cuts that edge, and the race
resolves to $B$. Occurrence-level stamping needs addresses into the term,
which the probes $\mathsf{Get}_{C,i}$ supply, so positional intervention
is a white-box operation. A \emph{full intervention} $\tau$ resolves
every race. Under a full intervention a flat process runs determinately,
up to the interleaving of independent steps.

\paragraph{Control.} Classical rho is the case $u = x = \quo{\nil}$ of the
graded comm \cite{FuzzyRho}. A copy sent at unit grade is the untreated
baseline.

\paragraph{Dose.} The order axiom of \cite{FuzzyRho} --- a larger match
grade transmits at least as much information --- is, for partial or soft
interventions, a monotone dose--response law. That law is what licenses
inference from graded treatments.

\paragraph{Readout.} The attenuator $u \cmp x = \quo{(\dr{u} \wedge
\dr{x})}$ can collapse distinct channels onto one shared stamped name.
Collapsing the success channels of a pool of copies onto one name gives a
single readout that reports whether \emph{any} member succeeded. A
negative pooled readout certifies failure for every member, provided the
collapse sends nothing else to that name. Pooling saves readouts, not
runs: each copy still realises one resolution.

\begin{remark}[What failure now means]
Under a full intervention $\tau$, failure is conclusive \emph{for $\tau$}.
The controlled apparatus therefore converts demonic scheduling into the
experimenter's choice. Success of $P$ is the existence of a $\tau$ under
which the copy succeeds --- an $\NPclass$ witness, now chosen rather than
granted. Failure of $P$ is failure under every $\tau$. The $\coNP$ half
becomes a problem of experimental design: can polynomially many
interventions, possibly pooled and graded, \emph{cover} the space of
resolutions?
\end{remark}

% =====================================================================
\section{Experimental regimes}\label{sec:regimes}
% =====================================================================

\begin{definition}[Regime]\label{def:regime}
An \emph{experimental regime} is a tuple $\Reg = (A, I, O, c)$ where
\begin{itemize}[leftmargin=2.2em,nosep]
\item $A$ is an access tier (Definition~\ref{def:tiers});
\item $I$ is an intervention language: a set of partial resolutions the
experimenter may impose, each realised by an experimental grading;
\item $O$ is a readout discipline: single, pooled or graded;
\item $c$ is the cost model, phlogiston satisfying Hypotheses~\ref{hyp:sim}
and~\ref{hyp:acct}.
\end{itemize}
A \emph{copy} is released by a graded communication that applies its
signature and its intervention; we identify the copy with that
communication. An \emph{experiment} is a finite pool of copies, run
against test contexts allowed by $A$, with a readout in $O$. $\Reg$ is
\emph{admissible} if, on every flat input, every experiment has an
outcome independent of the scheduler. $\enc{\Phi}(x)$ \emph{exploits}
$\Reg$ for $\ABS$ if, for every flat $P$, the configuration
$\enc{\Phi}(x) \mid \out{x}{P}$ performs only experiments of $\Reg$, runs
at phlogiston polynomial in $|P|$, and reports, independently of the
scheduler, whether $P$ can succeed. (Given this last clause,
admissibility is needed only of the experiments actually performed;
regimes whose interventions are partial are rarely admissible, and
Theorem~\ref{thm:blackbox} does not assume admissibility.)
\end{definition}

\begin{lemma}[Copies cost]\label{lem:copies}
Under Hypothesis~\ref{hyp:acct}, a run of phlogiston $c$ releases at most
$c$ copies.
\end{lemma}

\begin{proof}
Each copy is its own graded communication, charged at least $1$.
\end{proof}

\begin{remark}
Unstamped duplicates are cheaper. One communication can deliver $\quo{P}$
to a body containing many drops. But duplicates share their channels, so
they are neither isolated nor attributable, and none of them carries an
intervention. In the black-box model below an unintervened duplicate is
resolved by the scheduler, so it contributes nothing that an adversary
cannot neutralise.
\end{remark}

\subsection{Black-box regimes cannot certify universal negatives}

\begin{definition}[Black-box regime]\label{def:blackbox}
$\Reg$ is \emph{black-box} if its tier is (a) or (b) and the outcome of
each experiment on $P$ is determined by the experiment together with the
success or failure, at the stamped public interface, of each copy under
its intervention. Interventions are drawn from a public intervention
language over declared race ports. These are granted as a primitive,
which is more than a black-box experimenter could otherwise do, since
positional intervention needs addresses; granting them only strengthens
the lower bound below. Races left unresolved by an
intervention are resolved by a scheduler outside the experimenter's
control, and correctness is required under every scheduler.
\end{definition}

This definition is generous to the experimenter. It allows full
interventions without any inspection, and it allows pooled readouts and
any graded readout that is a function of the per-copy outcomes. It
excludes graded readouts that aggregate over the resolutions of a single
copy rather than taking the one the scheduler chose; those are the
subject of Conjecture~\ref{conj:algebrise}.

\begin{theorem}[Black-box lower bound]\label{thm:blackbox}
For each $d$ there are flat processes $P_\bot$ and $P_\sigma$ for
$\sigma \in \{0,1\}^d$, all of size $O(d \log d)$, such that $P_\bot$ cannot
succeed, each $P_\sigma$ can, and every $\enc{\Phi}(x)$ that, using a
black-box regime, reports correctly and independently of the scheduler
on all of them releases at least $2^d$ intervened copies on input
$P_\bot$. Under Hypothesis~\ref{hyp:acct} its phlogiston on $P_\bot$ is
at least $2^d$.
\end{theorem}

\begin{proof}
Let $\mathit{Race}_i = \out{b_i}{\nil} \mid \inp{\wild}{b_i}{\out{e_i^0}{\nil}}
\mid \inp{\wild}{b_i}{\out{e_i^1}{\nil}}$, with race port $r_i$ naming the
race on $b_i$. Put
\[
\begin{aligned}
P_\sigma &= \prod_{i \le d} \mathit{Race}_i \ \mid\
  \mathsf{for}(\wild \leftarrow e_1^{\sigma_1} \,\&\, \cdots \,\&\,
  \wild \leftarrow e_d^{\sigma_d})\, \out{\ok}{\nil},\\
P_\bot &= \prod_{i \le d} \mathit{Race}_i \ \mid\
  \mathsf{for}(\wild \leftarrow e_1^0 \,\&\, \cdots \,\&\,
  \wild \leftarrow e_d^0 \,\&\, \wild \leftarrow z)\, \out{\ok}{\nil}
\end{aligned}
\]
with $z$ a name nothing sends on, and all names pairwise distinct. (With
names as quoted numerals the size is $O(d \log d)$.) Then $P_\sigma$
succeeds exactly under the resolution $\sigma$, and $P_\bot$ never
succeeds.

Consider a copy with intervention $\tau$ (a partial assignment of race
outcomes). On $P_\bot$ it fails. On $P_\sigma$ it fails unless $\tau$ is
the full assignment $\sigma$: if $\tau$ disagrees with $\sigma$ somewhere,
the final join cannot fire; if $\tau$ leaves some race unresolved, a
scheduler that resolves it against $\sigma$ makes the copy fail. Fix such
a scheduler. By Definition~\ref{def:blackbox}, the outcome of every
experiment on $P_\sigma$ then coincides with its outcome on $P_\bot$,
unless the experiment contains a copy with full intervention $\sigma$.

Run $\enc{\Phi}(x)$ on $P_\bot$ under the same scheduler. It is correct,
so it reports failure. Let $T$ be the set of full interventions it ever applied
to a copy. If some $\sigma \notin T$, then on $P_\sigma$, under the
scheduler above, every experiment returns what it returned on $P_\bot$.
By induction on the run, $\enc{\Phi}(x)$ makes the same adaptive choices
and also reports failure, which is wrong. Hence $|T| = 2^d$. So
$\enc{\Phi}(x)$ released at least $2^d$ intervened copies on $P_\bot$, and
by Lemma~\ref{lem:copies} its phlogiston is at least $2^d$.
\end{proof}

\begin{corollary}\label{cor:nobb}
No black-box regime, admissible or not, is exploitable for $\ABS$. In
particular a linear regime, which admits one copy, is not.
\end{corollary}

\begin{remark}[Why this says nothing about $\Pclass$ versus $\NPclass$]
\label{rem:oracle}
In a white-box regime the family is trivial. $\enc{\Phi}$ reads the final
join and sees either the conjunct on $z$ or the hidden $\sigma$. The lower
bound is purely about access. It is a query lower bound for unstructured
search, the in-calculus analog of the \emph{separating} half of the
Baker--Gill--Solovay theorem \cite{BGS75}: a black-box experimenter
computes with $P$ as an oracle, and an unconditional separation is
available because the oracle hides its structure. We give no analog of
the \emph{collapsing} half, so the correspondence with the relativisation
barrier is an analogy, not a theorem.
\end{remark}

\begin{remark}[Graded and quantum readouts]
If a graded regime allowed amplitude-weighted superpositions of
interventions, as in the complex-weighted GSLTs, we would expect the bound
to weaken to $\Omega(2^{d/2})$ and no further, by the
Bennett--Bernstein--Brassard--Vazirani argument \cite{BBBV97}. Making this
precise requires a quantum semantics of the graded comm.
\end{remark}

\subsection{White-box regimes are as hard as \texorpdfstring{$\Pclass$}{P} versus \texorpdfstring{$\NPclass$}{NP}}

\begin{theorem}\label{thm:whitebox}
Assume Hypotheses~\ref{hyp:sim} and~\ref{hyp:acct}, and that reading a
term of size $n$ by probes costs phlogiston polynomial in $n$. Then
$\Pclass = \NPclass$ if and only if some admissible regime is exploitable
for $\ABS$. When it is, a white-box regime is.
\end{theorem}

\begin{proof}
($\Rightarrow$) If $\Pclass = \NPclass$ then $\ABS \in \Pclass$ by
Proposition~\ref{prop:abscomplete}, decided by some machine $M$. In a
white-box regime $\enc{\Phi}(x)$ receives $\quo{P}$, reads the quoted code
by probes into a data representation at polynomial cost, and runs $R_M$
on it (Hypothesis~\ref{hyp:sim}). No copy of $P$ ever runs, so its report
does not depend on the scheduler.

($\Leftarrow$) Suppose $\enc{\Phi}(x)$ exploits a regime $\Reg$. Its report
is independent of the scheduler, so simulating any one schedule computes
it, and the whole configuration runs at polynomial phlogiston. By Hypothesis~\ref{hyp:acct} a deterministic
machine simulates it in polynomial time, so $\ABS \in \Pclass$ and hence
$\coNP = \Pclass = \NPclass$.
\end{proof}

The direction ($\Leftarrow$) holds for every regime, black-box or
white-box, and uses admissibility only through the scheduler-independence
of the report. Theorem~\ref{thm:whitebox} is a restatement, not a new
technique: its content is in combination with Theorem~\ref{thm:blackbox}.

\begin{corollary}[The restatement]\label{cor:restatement}
Under the hypotheses of Theorem~\ref{thm:whitebox},
\[
\Pclass \neq \NPclass
\quad\Longleftrightarrow\quad
\begin{array}{l}
\text{no admissible regime lets } \enc{\Phi}(x) \text{ certify, at}\\
\text{polynomial cost, that a flat process cannot succeed.}
\end{array}
\]
Black-box regimes provably cannot (Theorem~\ref{thm:blackbox}), so the
remaining content of $\Pclass \neq \NPclass$, in this formulation, is that
white-box access affords no exploit. A proof of $\Pclass = \NPclass$ would
exhibit a laboratory.
\end{corollary}

% =====================================================================
\section{The barriers, located}\label{sec:barriers}
% =====================================================================

The correspondences in this section are analogies. We record them because
they say where in the apparatus each barrier would bite.

\paragraph{Relativisation and black-box access.} An argument that applies
uniformly to black-box regimes treats the input as an oracle.
Theorem~\ref{thm:blackbox} and Remark~\ref{rem:oracle} show that such
arguments can separate unconditionally while saying nothing about the
unrelativised question. This mirrors the barrier of \cite{BGS75}. A full
analog would also need a regime playing the role of the collapsing
oracle; a regime with a primitive experiment answering a
$\mathbf{PSPACE}$-complete question would do, but it is not admissible in
our sense, since that experiment does not cost polynomial phlogiston.

\paragraph{Algebrisation and graded aggregate readout.} A graded readout
of the kind excluded from Definition~\ref{def:blackbox} returns an
aggregate --- a weighted sum over the resolutions of a copy --- rather
than the one outcome the scheduler chose. Heuristically this resembles
querying an extension of the oracle beyond its truth table, the access
model of the algebrisation barrier \cite{AaronsonWigderson09}, though a
weighted sum is not a low-degree extension over a larger field. We
conjecture that graded black-box regimes form a second tier that is
unconditionally separated and equally uninformative
(Conjecture~\ref{conj:algebrise}).

\paragraph{Natural proofs.} An argument that certifies hardness through a
large, constructive property of a process's full behaviour --- its table
of outcomes over all interventions --- is the analog of a natural proof
\cite{RazborovRudich97}. Making this precise is open.

\begin{conjecture}\label{conj:algebrise}
Theorem~\ref{thm:blackbox} holds, with a bound exponential in $d$, for
black-box regimes extended by graded readouts that aggregate over the
resolutions of each copy, provided grades are computed from interface
names.
\end{conjecture}

% =====================================================================
\section{Discussion}\label{sec:discussion}
% =====================================================================

\paragraph{Popper's asymmetry, with a price.} ``$P$ can succeed'' is
existential. One experiment verifies it, and no number of failed
experiments falsifies it until every resolution has been tried. ``$P$
cannot succeed'' is universal. One experiment falsifies it, and verifying
it means exhausting the resolutions \cite{Popper59}.
Corollary~\ref{cor:restatement} puts a price on the asymmetry. Black-box
laboratories pay exponentially, provably. Whether a white-box laboratory
can pay polynomially is $\Pclass$ versus $\NPclass$.

\paragraph{The mortal scientist.} A scientist modelled as a
cost-accounted rho computation, observing populations through graded
where-clauses, works in a black-box regime equipped with a distribution
over schedules. Its evidence of absence is statistical --- a tail bound,
not a certificate. The randomised analog of
Corollary~\ref{cor:restatement} would compare such regimes with
$\mathbf{BPP}$ or $\mathbf{AM}$ rather than $\Pclass$, and may be the more
tractable version inside cost-accounted rho, because the graded structure
already measures how much a failure tells.

\paragraph{Automata learning.} Repeated black-box experiments with
steerable choice are membership queries. Determinate systems are
learnable from membership and equivalence queries in polynomial time
\cite{Angluin87}. Under cryptographic assumptions, membership queries do
not help for nondeterministic automata \cite{AngluinKharitonov95}. This
matches the determinate/nondeterminate split of
Proposition~\ref{prop:determinate} and Theorem~\ref{thm:blackbox}, though
the correspondence has not been worked out.

\paragraph{Independence, revisited.} If $\Pclass \neq \NPclass$ were
independent of ZFC, ZFC could not rule out a laboratory. In some model
with nonstandard integers, a nonstandard $\enc{\Phi}$ would exploit a
regime whose cost bound is a nonstandard polynomial. The theorem of
\cite{BenDavidHalevi92} concerns a different theory, PA with all true
$\Pi_1$ sentences, and does not apply to ZFC directly. It does indicate
how close to real such laboratories would have to be if independence held
for theories of that kind. The
Boolean-valued models of \cite{Krajicek11} have random variables as
their integers. Grades valued in such an algebra are a natural candidate
for the fuzzy rho of those models, and connecting the two is open.

% =====================================================================
\section{Open problems}\label{sec:open}
% =====================================================================

\begin{enumerate}[label=\arabic*.,leftmargin=2.2em]
\item \textbf{Compositional normalisers.} Prove
Conjecture~\ref{conj:compositional}. It is the analog of a monotone
lower bound, and the most promising source of a provable restricted
separation.

\item \textbf{Equivariant stamping.} Realise (E3) for general rho
processes by a namespace-relative evaluation in which every name the copy
synthesises is stamped, or show that no rho context achieves it.

\item \textbf{Drop-only congruence.} Prove
Conjecture~\ref{conj:dropOnly}, making tier (b) an admissible observer
class in the sense of \cite{BIHOL}.

\item \textbf{Graded regimes.} Prove Conjecture~\ref{conj:algebrise}, and
identify the quantum analog with its $\Omega(2^{d/2})$ bound.

\item \textbf{Statistical regimes.} State and prove the randomised
analog of Corollary~\ref{cor:restatement}.

\item \textbf{Minimal labels.} Replace enabling contexts by minimal ones,
following \cite{LeiferMilner00,BIHOL}, and check that
Proposition~\ref{prop:adequacy} survives.

\item \textbf{Cost hypotheses.} Verify Hypotheses~\ref{hyp:sim}
and~\ref{hyp:acct} for cost-accounted rho as implemented, and determine
the cost of probes in the observer extension.

\item \textbf{The intermediate tier.} Characterise regimes with
positional interventions but no full reading of the term. Is there an
access discipline strictly between (b) and (c) that is still provably
separated?
\end{enumerate}

% =====================================================================
\begin{thebibliography}{99}
% =====================================================================

\bibitem{Aaronson03}
S.~Aaronson.
\newblock Is $P$ versus $NP$ formally independent?
\newblock \emph{Bulletin of the EATCS}, 81:109--136, 2003.

\bibitem{AaronsonWigderson09}
S.~Aaronson and A.~Wigderson.
\newblock Algebrization: a new barrier in complexity theory.
\newblock \emph{ACM Transactions on Computation Theory}, 1(1), 2009.

\bibitem{Angluin87}
D.~Angluin.
\newblock Learning regular sets from queries and counterexamples.
\newblock \emph{Information and Computation}, 75(2):87--106, 1987.

\bibitem{AngluinKharitonov95}
D.~Angluin and M.~Kharitonov.
\newblock When won't membership queries help?
\newblock \emph{Journal of Computer and System Sciences}, 50(2):336--355,
1995.

\bibitem{BGS75}
T.~Baker, J.~Gill, and R.~Solovay.
\newblock Relativizations of the $P =? NP$ question.
\newblock \emph{SIAM Journal on Computing}, 4(4):431--442, 1975.

\bibitem{BGS92}
J.~L. Balc\'azar, J.~Gabarr\'o, and M.~S\'antha.
\newblock Deciding bisimilarity is $P$-complete.
\newblock \emph{Formal Aspects of Computing}, 4(6A):638--648, 1992.

\bibitem{BenDavidHalevi92}
S.~Ben-David and S.~Halevi.
\newblock On the independence of $P$ versus $NP$.
\newblock Technical Report 714, Technion, 1992.

\bibitem{BBBV97}
C.~H. Bennett, E.~Bernstein, G.~Brassard, and U.~Vazirani.
\newblock Strengths and weaknesses of quantum computing.
\newblock \emph{SIAM Journal on Computing}, 26(5):1510--1523, 1997.

\bibitem{FuzzyRho}
Fuzzy rho-calculus: a quantale-graded interaction, and its lift to
ci\,GSLTs with binding.
\newblock Working note, F1R3FLY.io publications, \texttt{fuzzy-rho}.

\bibitem{Krajicek11}
J.~Kraj\'{\i}\v{c}ek.
\newblock \emph{Forcing with Random Variables and Proof Complexity}.
\newblock Volume 382 of LMS Lecture Note Series. Cambridge University
Press, 2011.

\bibitem{LeiferMilner00}
J.~J. Leifer and R.~Milner.
\newblock Deriving bisimulation congruences for reactive systems.
\newblock In \emph{CONCUR 2000}, volume 1877 of \emph{LNCS}, pages
243--258. Springer, 2000.

\bibitem{NWNS}
L.~G. Meredith.
\newblock Namespaces without a name server: freshness by communication,
and a fully abstract encoding of the $\pi$-calculus in the rho calculus.
\newblock Draft 3, F1R3FLY.io publications, \texttt{pi2rhobugfix}, 2026.

\bibitem{ObserverExt}
L.~G. Meredith.
\newblock Every constructor an interaction: observer extensions, and the
discriminating power of spatial formulae in interactive graph-structured
lambda theories.
\newblock Draft, F1R3FLY.io publications,
\texttt{spatial-formulae-discrimination-power}, 2026.

\bibitem{MeredithRadestock05}
L.~G. Meredith and M.~Radestock.
\newblock A reflective higher-order calculus.
\newblock \emph{Electronic Notes in Theoretical Computer Science},
141(5):49--67, 2005.

\bibitem{PaigeTarjan87}
R.~Paige and R.~E. Tarjan.
\newblock Three partition refinement algorithms.
\newblock \emph{SIAM Journal on Computing}, 16(6):973--989, 1987.

\bibitem{MeyerFischer71}
A.~R. Meyer and M.~J. Fischer.
\newblock Economy of description by automata, grammars, and formal
systems.
\newblock In \emph{12th Annual Symposium on Switching and Automata
Theory}, pages 188--191. IEEE, 1971.

\bibitem{Moore71}
F.~R. Moore.
\newblock On the bounds for state-set size in the proofs of equivalence
between deterministic, nondeterministic, and two-way finite automata.
\newblock \emph{IEEE Transactions on Computers}, C-20(10):1211--1214,
1971.

\bibitem{Pearl09}
J.~Pearl.
\newblock \emph{Causality}.
\newblock Cambridge University Press, second edition, 2009.

\bibitem{Popper59}
K.~R. Popper.
\newblock \emph{The Logic of Scientific Discovery}.
\newblock Hutchinson, 1959.

\bibitem{Razborov95}
A.~A. Razborov.
\newblock Unprovability of lower bounds on circuit size in certain
fragments of bounded arithmetic.
\newblock \emph{Izvestiya: Mathematics}, 59(1):205--227, 1995.

\bibitem{RazborovRudich97}
A.~A. Razborov and S.~Rudich.
\newblock Natural proofs.
\newblock \emph{Journal of Computer and System Sciences}, 55(1):24--35,
1997.

\bibitem{Shoenfield61}
J.~R. Shoenfield.
\newblock The problem of predicativity.
\newblock In \emph{Essays on the Foundations of Mathematics}, pages
132--139. Magnes Press, 1961.

\bibitem{StockmeyerMeyer73}
L.~J. Stockmeyer and A.~R. Meyer.
\newblock Word problems requiring exponential time.
\newblock In \emph{STOC 1973}, pages 1--9. ACM, 1973.

\bibitem{BIHOL}
C.~B. Wells and M.~Stay.
\newblock Behavior in higher-order languages.
\newblock In preparation.

\end{thebibliography}

\appendix

% =====================================================================
\section{Status of results}\label{app:status}
% =====================================================================

\begin{center}
\begin{tabular}{@{}llp{0.46\textwidth}@{}}
\toprule
Result & Status & Note \\
\midrule
Prop.~\ref{prop:adequacy} & sketched & administrative phases as in \cite{NWNS} \\
Thm.~\ref{thm:known} & known & (b) via \cite{StockmeyerMeyer73}; reduction sketched \\
Faithfulness remark after Prop.~\ref{prop:determinate} & proved & \\
Prop.~\ref{prop:abscomplete} & proved & \\
Thm.~\ref{thm:normaliser} & proved & \\
Prop.~\ref{prop:determinate} & proved & unconditional \\
Conj.~\ref{conj:compositional} & conjecture & open problem 1 \\
Lem.~\ref{lem:uses} & proved & syntactic \\
Ex.~\ref{ex:linear} & illustrative & formal content in Thm.~\ref{thm:blackbox} \\
Conj.~\ref{conj:dropOnly} & conjecture & open problem 3 \\
Lem.~\ref{lem:injective}, \ref{lem:disjoint} & proved & \\
Ex.~\ref{ex:nonequivariant} & proved & \\
Prop.~\ref{prop:fidelity} & proved & given (E1)--(E3) \\
Lem.~\ref{lem:copies} & proved & relative to Hyp.~\ref{hyp:acct} and Def.~\ref{def:regime} \\
Thm.~\ref{thm:blackbox}, Cor.~\ref{cor:nobb} & proved & in the model of Def.~\ref{def:blackbox}; copy count unconditional, cost via Hyp.~\ref{hyp:acct} \\
Thm.~\ref{thm:whitebox}, Cor.~\ref{cor:restatement} & proved & relative to Hyp.~\ref{hyp:sim}, \ref{hyp:acct} and probe cost; a restatement \\
\S\ref{sec:barriers} & analogy & not theorems \\
Conj.~\ref{conj:algebrise} & conjecture & open problem 4 \\
\bottomrule
\end{tabular}
\end{center}

\end{document}
